\documentclass[11pt]{article}
\usepackage{acl}
\usepackage{times}
\usepackage{latexsym}
\usepackage[T1]{fontenc}
\usepackage[utf8]{inputenc}
\usepackage{microtype}
\usepackage{inconsolata}
\usepackage{graphicx}

\title{HedgeBERT: Hedge-Aware Sentiment Classification via CLS\@ Feature Injection}

\author{
  Minh Tran \\
  Wichita State University \\
  \texttt{K699E569@wichita.edu}
  \And
  An Nguyen \\
  Wichita State University \\
  \texttt{R374G648@wichita.edu}
}

\begin{document}
\maketitle

\begin{abstract}
% Write your abstract here (2-4 sentences summarizing the problem, your approach, and expected contribution)
Sentiment analysis models struggle with classifying hedged language, as it can be less direct and ambiguous. 
Hedge Language is also common in many real world scenarios such as social media posts or product reviews.
Mistaking a negative review with hedged language for a positive one on a product can lead to a company not evaluating their product properly.
We propose HedgeBERT which is a modified BERT model which encodes a hedge intensity score into the CLS token to help BERT better make decisions on hedged languages.
We expect HedgeBERT to demonstrate measurable accuracy improvement on hedged consumer sentiment without degrading performance on direct sentiment.
\end{abstract}

\section{Introduction}
% What is the problem? Why does it matter? 
% Briefly introduce hedged language and why sentiment classifiers fail on it.
% End with a summary of your contribution (HedgeBERT).

\section{Related Work}
% Summarize existing work on:
% 1. Sentiment analysis (Pang & Lee, BERT-based models)
% 2. Hedge detection (BioScope, SFU corpus)
% 3. What gap exists (nobody has used hedging to improve sentiment classification)

\section{Research Question}
% Clearly state your research question:
% Can explicitly encoding hedge intensity as a feature improve 
% the accuracy of BERT-based sentiment classifiers on hedged text?

\section{Proposed Methodology}
% Describe your three-part approach:
% 1. Hedged-Sentiment Benchmark Dataset (~250 examples from Yelp/Amazon)
% 2. Hedge Intensity Scorer (rule-based, outputs 0.0 to 1.0)
% 3. CLS Feature Injection (concatenate hedge score to [CLS] before classification)

\section{Evaluation Plan}
% How will you measure success?
% Compare HedgeBERT vs standard DistilBERT on hedged vs direct sentences
% Metrics: Accuracy, F1 Score

\section{Timeline}
% Rough timeline of milestones:
% Week 1-2: Dataset curation, hedge scorer
% Week 3-4: Baseline model setup, CLS injection
% Week 5-6: Fine-tuning, progress presentations (April 8 & 13)
% Week 7: Evaluation
% Week 8: Final report and presentation

% Bibliography
%\bibliography{custom}
% uncomment once some things are in the bib i think

\end{document}